% !TeX root = ../../Book4.tex
% 纲要标题：### 13.1 余极限的正合性
% 纲要位置：计划纲要.md，第 2799 行。

\section{余极限的正合性}
\label{b4:sec:1301}

% 纲要内容（仅作写作计划，尚非教材正文）：
%
% * 模范畴中的滤过余极限
% * 一般余极限并非都正合
% * 积与直和的正合性依赖所在范畴
%

\subsection{模范畴中的滤过余极限}
\label{b4:subsec:1301:01}

滤过余极限允许把有限次运算放到同一个阶段进行，而逆极限要求一整族元素同时相容。这两种条件的差别，决定了它们的正合性质。本章在模范畴中讨论这些问题；若换成其他交换范畴，还须重新检查所用的无限积是否正合。

\ProseParagraphBreak
\cref{b4:prop:filtered-colimits-exact}已证明：对小滤过范畴上的模图表，逐对象的短正合列在取余极限后仍正合。证明的关键是\cref{b4:prop:filtered-module-elements}的单阶段判据：一个元素在余极限中为零，当且仅当它在某个更后阶段已经变成零。下面将此结果用于复形，而不重作余极限的构造。

\begin{proposition}\label{b4:prop:filtered-colimits-cohomology}
设 \(R\) 为环，\(\mathcal I\) 为小滤过范畴，\(C_i\) 为左 \(R\)-模上链复形组成的图表。则对每个整数 \(q\)，有自然同构
\[
\varinjlim_{i\in\mathcal I}H^q(C_i)
\simeq H^q\bigl(\varinjlim_{i\in\mathcal I}C_i\bigr).
\]
右侧余极限逐次数形成，不要求复形有界。
\end{proposition}
\begin{proof}
把复形的微分拆成短正合列
\[
0\to Z^q(C_i)\to C_i^q\to B^{q+1}(C_i)\to0,\qquad
0\to B^q(C_i)\to Z^q(C_i)\to H^q(C_i)\to0.
\]
余极限保持这两列正合，也保持 \(B^{q+1}(C_i)\to C_i^{q+1}\) 的单射性。因此取余极限后的圈与边界分别是原圈与边界的余极限，它们之商正是 \(\varinjlim H^q(C_i)\)。所有识别使用自然包含及商映射，故自然。每次仅涉及两个相邻微分，与是否有界无关。
\end{proof}

例如，一个复形写成递增子复形的并时，任一同调类可以在某个阶段表示；若它最终变成边界，也在某个阶段已经变成边界。这不是关于极限的形式类比，而是有限代表元判据的实际结果。

\subsection{一般余极限的非正合性}
\label{b4:subsec:1301:02}

一般余极限是常值图表函子的左伴随，所以保持余核，因而右正合；要同时保持核，指标范畴的条件就有作用。

\begin{example}\label{b4:exm:parallel-colimit-not-exact}
令 \(\mathcal I\) 只有对象 \(0,1\) 及两支不同箭头 \(s,t:0\to1\)。交换群图表 \(X\) 的余极限是 \(\operatorname{coker}(X(s)-X(t))\)。定义三个图表
\[
\begin{array}{c|cc|cc}
 &X(0)&X(1)&X(s)&X(t)\\\hline
A&0&\mathbb Z&0&0\\
B&\mathbb Z&\mathbb Z&\times2&0\\
C&\mathbb Z&0&0&0
\end{array}
\]
在对象 \(0\) 处取 \(0\to0\to\mathbb Z\xrightarrow{1}\mathbb Z\to0\)，在对象 \(1\) 处取 \(0\to\mathbb Z\xrightarrow{1}\mathbb Z\to0\to0\)，得到图表短正合列 \(0\to A\to B\to C\to0\)。它的余极限列为
\[
\mathbb Z\longrightarrow\mathbb Z/2\longrightarrow0.
\]
首箭头不是单射。表中的两个平行箭头不能被任何更后箭头合并；这个指标范畴不滤过。
\end{example}

这个例子的关系 \(2\mathbb Z\) 来自图表的箭头。逐对象检查单射，并不能阻止余极限在目标中额外施加这些关系。因而“每个阶段正合”只有结合正确的指标条件，才足以推出余极限正合。

\subsection{积与直和的正合性条件}
\label{b4:subsec:1301:03}

对左 \(R\)-模，任意一族短正合列取直和仍正合：一项元素只有有限支集，可以在这有限多个位置选取原像。取直积也仍正合，但满射性的证明需要对全部坐标同时选原像，使用集合论选择。核与像的其余核对均逐坐标进行。

\begin{proposition}\label{b4:prop:products-cohomology}
设 \(R\) 为环，\((C_j)_{j\in J}\) 为集合指标的左 \(R\)-模上链复形族。则对每个 \(q\)，有自然同构
\[
H^q\bigl(\prod_j C_j\bigr)\simeq\prod_jH^q(C_j),\qquad
H^q\bigl(\bigoplus_j C_j\bigr)\simeq\bigoplus_jH^q(C_j).
\]
\end{proposition}
\begin{proof}
积复形的圈逐坐标给出。对一族边界，每个坐标选一个微分原像，就得到积中的原像，所以积的边界也是边界的积。对 \(0\to B^q\to Z^q\to H^q\to0\) 取积并使用满射性，便得到第一式。直和时只需在有限个非零坐标上选择，完全相同的核、像、商计算给出第二式。
\end{proof}

在一般交换范畴中，存在积与积正合是两个不同要求。若范畴有小余积且余积保持短正合列，通常称其满足 \(\mathrm{AB4}\)；若有小积且积保持短正合列，称其满足 \(\mathrm{AB4}^*\)。若小滤过余极限存在且正合，则称其满足 \(\mathrm{AB5}\)。这些名称描述附加性质，不是交换范畴定义的组成部分。

\begin{example}\label{b4:exm:torsion-category-product}
所有挠交换群及其同态组成交换范畴 \(\mathcal T\)：核、余核和有限双积仍是挠群。该范畴内一族群的积，是通常直积的挠子群，因为从任何挠群发出的同态必然落入这个子群。

\ProseParagraphBreak
固定素数 \(p\)，对每个 \(n\ge1\) 取满射 \(\mathbb Z/p^n\to\mathbb Z/p\)。目标在 \(\mathcal T\) 中的积等于通常直积，其中 \((1,1,\ldots)\) 的阶为 \(p\)。它若有原像，原像的第 \(n\) 个坐标必须是模 \(p\) 等于 \(1\) 的元素，阶恰为 \(p^n\)。该元组没有有限阶，因而不属于源在 \(\mathcal T\) 中的积。这说明 \(\mathcal T\) 的积不保持满射。
\end{example}

后面乘积差映射的证明明确限定在模范畴中。更广版本至少需要相应的积正合性；不能仅凭“也是交换范畴”便移用。
